\subsection{Ariel Data Challenge 2025}

Mục đích của cuộc thi ADC2025 là dự đoán phần trăm mất mát ứng với từng bước sóng phổ của ánh sáng cùng với
độ không chắc chắn. Mất mát này xảy ra khi ánh sáng từ một ngôi sao chủ đi qua bầu khí quyển của một hành 
tinh quay quanh nó. Khi ánh sáng đi qua bầu khí quyển, một số bước sóng sẽ bị hấp thụ hoặc tán xạ bởi các
phân tử trong khí quyển, dẫn đến việc giảm cường độ ánh sáng \textit{đối với các bước sóng cụ thể}. Bằng cách
phân tích sự mất mát đó, ta có thể xác định thành phần hóa học của bầu khí quyển hành tinh, từ đó suy ra các
đặc tính như nhiệt độ, áp suất, thành phần khí hay quan trọng hơn là xác định sự sống trên hành tinh đó. 

Để những phương pháp trong cuộc thi có thể được áp dụng rộng rãi, dữ liệu mô phỏng trong cuộc thi được cung cấp
rất sát với thực tế, giúp các nhà nghiên cứu và phát triển có thể kiểm tra và đánh giá hiệu quả của các phương pháp
mới trong một môi trường gần với thực tế nhất có thể.